\documentclass[a4paper,8pt]{extarticle}

% Packages
\usepackage[T1]{fontenc}
\usepackage[utf8]{inputenc}
\usepackage{geometry}
\usepackage{hyperref}
\usepackage{fontawesome5}
\usepackage{titlesec}
\usepackage{enumitem}
\usepackage{xcolor}
\usepackage{tabularx}
\usepackage{multirow}
\usepackage{tikz}
\usepackage{lmodern}  

% Define custom colors
\definecolor{primary}{RGB}{11, 25, 47}
\definecolor{darkgray}{RGB}{50, 50, 50}
\definecolor{lightgray}{RGB}{245, 245, 245}
\definecolor{linkcolor}{RGB}{70, 130, 210}

% Configure hyperlinks - updated for better clickability
\hypersetup{
  colorlinks=true,
  linkcolor=linkcolor,
  filecolor=linkcolor,
  urlcolor=linkcolor,
  citecolor=linkcolor,
  pdftitle={Lucian Medrihan CV},
  pdfauthor={Lucian Medrihan},
  pdfborder={0 0 1},
  breaklinks=true,
  bookmarksopen=true,
  bookmarksnumbered=true
}



% Page geometry
\geometry{
  top=0.6cm,
  bottom=0.6cm,
  left=1cm,
  right=1cm
}

% Remove paragraph indentation
\setlength{\parindent}{0pt}

% Section formatting
\titleformat{\section}
  {\large\bfseries\color{primary}}
  {}{0em}{}[\titlerule]
\titlespacing{\section}{0pt}{6pt}{3pt}

% Skills formatting
\newcommand{\skill}[1]{
  \colorbox{lightgray}{\footnotesize #1}
}

% Timeline dot
\newcommand{\timelinedot}{
  \tikz\draw[primary, fill=primary] (0,0) circle (2pt);
}

% Experience item
\newcommand{\experienceitem}[5]{
  \timelinedot~\textbf{#1}\\
  \textit{#2} \hfill \textit{#3}\\
  #4\\
  \begin{itemize}[leftmargin=1em,noitemsep,topsep=0pt,parsep=0pt]
    #5
  \end{itemize}
}


\newcommand{\projectitem}[5]{
  \textbf{#1} \hfill \textit{#2}\\
  #3 \hfill 
  \ifnum\pdfstrcmp{#4}{null}=0\else\href{#4}{\faGlobe}\fi
  \ifnum\pdfstrcmp{#5}{null}=0\else~\href{#5}{\faGithub}\fi\\
  \begin{itemize}[leftmargin=1em,noitemsep,topsep=0pt,parsep=0pt]
}

% Rating stars
\newcommand{\rating}[1]{
  \foreach \n in {1,...,5}{
    \ifnum\n>#1
      \textcolor{lightgray}{$\star$}
    \else
      \textcolor{primary}{$\star$}
    \fi
  }
}

\begin{document}

% Start two-column layout
\begin{minipage}[t]{0.33\textwidth}
  % Left column content
  \begin{center}
    {\fontsize{21}{20}\selectfont\textbf{\textcolor{primary}{Lucian Medrihan}}}\\
    \vspace{0.2cm}
    {\large Ingegnere Software}
  \end{center}
  
  \section{Chi Sono}
  Ingegnere con oltre 5 anni di esperienza nella realizzazione di sistemi in produzione per istituzioni pubbliche. Mi dedico a codice che viene effettivamente rilasciato e scala, da pipeline di immagini satellitari a giochi mobili offline. Ho contribuito a progetti per il SIAN e l'ESA dimostrando forti capacità di problem-solving.
  \section{Contatti}

\begin{tabular}{@{}p{0.1\textwidth}p{0.9\textwidth}@{}}
    \faPhone & \href{tel:+393272898921}{+39 3272898921} \\
    \faEnvelope & \href{mailto:medrihanlucian@gmail.com}{medrihanlucian@gmail.com} \\
    \faGlobe & \href{https://1vcian.github.io/Portfolio/}{1vcian.github.io/Portfolio} \\
    \faLinkedin & \href{https://linkedin.com/in/lume}{linkedin.com/in/lume} \\
    \faGithub & \href{https://github.com/1vcian}{github.com/1vcian} \\
    \faMapMarker & Roma, Italia \\
  \end{tabular}
  
  \section{Competenze}
  \textbf{Competenze Tecniche:} Esperto in \textcolor{linkcolor}{\textbf{Python}}, JavaScript e Node.js, con vasta esperienza in \textcolor{linkcolor}{\textbf{React.js}}, \textcolor{linkcolor}{\textbf{Angular.js}}, Flask e FastAPI. Specializzato nello \textcolor{linkcolor}{\textbf{sviluppo full-stack}} con tecnologie web moderne, completato da solide competenze in \textcolor{linkcolor}{\textbf{Docker}}, \textcolor{linkcolor}{\textbf{AWS}} e soluzioni cloud.\\
  
  \textbf{Competenze Aggiuntive:} \textcolor{linkcolor}{\textbf{GIS}}, \textcolor{linkcolor}{\textbf{Machine Learning}}, \textcolor{linkcolor}{\textbf{LLMs}}, AI-assisted development, Elaborazione Immagini, Tailwind CSS, Bootstrap, Openlayers, comandi Linux, \textcolor{linkcolor}{\textbf{Cybersecurity}}, Geopandas, JQuery, programmazione di rete.\\
  
  \textbf{Soft Skills:} Comprovata capacità di risolvere \textcolor{linkcolor}{\textbf{problemi complessi}}, gestire progetti in modo efficace e \textcolor{linkcolor}{\textbf{guidare team}}. Forte gestione del tempo, pensiero critico e attenzione ai dettagli.
  
  \section{Formazione}
  \timelinedot~\textbf{Laurea in Ingegneria Informatica e dell'Automazione}\\
  \textit{Università Sapienza di Roma}
  

  \section{Lingue}
  \begin{tabular}{@{}p{0.4\textwidth}p{0.6\textwidth}@{}}
    Italiano & \rating{5} \\
    Rumeno & \rating{5} \\
    Inglese & \rating{4} \\
  \end{tabular}

  \section{Riconoscimenti \& Attività}
  \begin{itemize}[leftmargin=1.5em]
    \item \textcolor{linkcolor}{\textbf{Bug Bounty Hunter}}: Trovato 2 RCE critiche (CVSS 10.0) su major ISP italiano via Cyberdart.eu (~\$5k bounty)
    \item \textcolor{linkcolor}{\textbf{Vincitore Google Workshop}} per UniverCity, presentato presso Google Tel Aviv
    \item Co-Fondatore di \textcolor{linkcolor}{\textbf{Hackappatoi}}, il team di ethical hacking e CTF di Sapienza
  \end{itemize}
  
\end{minipage}
\hfill
\begin{minipage}[t]{0.65\textwidth}
  % Right column content
  \section{Esperienze}
  
  \experienceitem{Ingegnere Software}{PR.EL.CART. EUROPE ENGINEERING S.R.L.}{Giu 2025 - Presente}{Roma, Italia}{
    \item Sviluppo in collaborazione con il Ministero dell'Agricoltura italiano per soluzioni software avanzate
    \item Gestione e mentoring di un team di 3 sviluppatori junior, guidandoli nell'architettura software
    \item Sviluppo indipendente dell'applicazione ufficiale \href{https://agea-vinitaly-2026.sian.it/}{AGEA Vinitaly 2026} utilizzando la tecnologia PMTiles
    \item Creazione di una pipeline AI su Google Cloud Vertex AI per generare foto satellitari a 1m da foto Sentinel a 10m, riducendo i costi da ~4/5M€ a ~500€/mese
  }
  
  \vspace{0.25cm}
  
  \experienceitem{Ingegnere Software}{Tiuke S.R.L.}{Ott 2019 - Giu 2025}{Roma, Italia}{
    \item Sviluppo di 5+ applicazioni web e mobile per il Ministero dell'Agricoltura, servendo oltre 10.000 utenti attivi
    \item Collaborato in un team di sviluppo ad alte prestazioni di 5 membri con pari potere decisionale, ottimizzando i rilasci e il coordinamento dei task
    \item Ottimizzato applicazioni legacy monolitiche in web app responsive, riducendo i tempi di caricamento del 40\%
    \item Creato strumenti di automazione per l'integrazione di layer cartografici, accelerando i tempi di sviluppo di 3x
    \item Sviluppato sistemi backend per la generazione di report, risparmiando circa 15 ore di lavoro manuale a settimana
    \item Implementato agenti LLM specializzati per l'estrazione di dati agricoli, raggiungendo il 95\% di accuratezza
    \item Migliorato l'addestramento di modelli di segmentazione geospaziale, incrementando la precisione del 20\%
    \item Scritto script batch per processare terabyte di dati, riducendo il tempo di esecuzione delle pipeline del 50\%
  }
  
  % Right column content
  \section{Progetti \hfill \normalfont\small\href{https://github.com/1vcian}{Vedi altri su GitHub}}
  
  \projectitem{Card Scanner [GitHub: $\star$ 58]}{Completato}{Sistema di riconoscimento carte Pokémon tramite AI con elaborazione client}{https://1vcian.github.io/Pokemon-TCGP-Card-Scanner/}{https://github.com/1vcian/Pokemon-TCGP-Card-Scanner}
    \item Costruito un sistema leggero (<10MB) di identificazione carte lato client che funziona con tutte le nuove carte senza richiedere riaddestramento
    \item Integrato con successo come scanner core MVP nell'open-source \href{https://github.com/marcelpanse/tcg-pocket-collection-tracker}{TCG Pocket Collection Tracker} ($\star$ 150+), migliorando significativamente l'esperienza utente
  \end{itemize}
    
  \vspace{0.25cm}
  
  \projectitem{EGNSS Capacitor}{Completato}{Plugin Capacitor universale per GNSS con supporto OSNMA}{https://1vcian.me/egnss-capacitor/}{https://github.com/1vcian/egnss-capacitor}
    \item Sviluppato un plugin multipiattaforma per garantire una precisione geografica assoluta in Europa sfruttando OSNMA ed EGNOS di Galileo
    \item Esposto un'API unificata (Android, iOS, Web) per acquisire coordinate ad altissima precisione da GNSS interno e antenne Bluetooth esterne
  \end{itemize}
    
  \vspace{0.25cm}


  \projectitem{Wolf \& Mafia}{Completato}{Gioco 100\% offline di Lupus in Tabula per Android e iOS}{https://wolfmafia.app/}{null}
    \item Sviluppato da zero un gioco completo di deduzione sociale con 36 ruoli unici e un sistema di narratore totalmente guidato
    \item Basato su un sistema di seed deterministico per gestire lo stato della partita completamente offline senza server o internet
    \item Implementato multiplayer locale basato su codice QR, permettendo agli utenti di giocare all'istante
  \end{itemize}
    
  \vspace{0.25cm}

  \projectitem{Strumenti per Gaming Community}{In Corso}{App ForgeMaster \& Eatventure Seed Predictor}{null}{null}
    \item Sviluppato un'app compagna auto-aggiornante per ForgeMaster (\href{https://1vcian.me/fm}{\faGlobe} / \href{https://github.com/1vcian/fm}{\faGithub})
    \item Creato un accurato predittore deterministico di drop RNG per Eatventure (\href{https://eatventure-loot-predictor.vercel.app/}{\faGlobe} / \href{https://github.com/1vcian/eatventure-vercel}{\faGithub})
  \end{itemize}
    
  \vspace{0.25cm}

  \projectitem{Smart House}{Completato}{Sistema domotico personalizzato basato su Arduino}{null}{https://github.com/LNPacio/ProgettoOS}
    \item Sviluppato un sistema domotico completo su Arduino MEGA2560 con accesso remoto sicuro via SSH
    \item Implementato controlli PWM per luci, ingressi sensori ADC e un'interfaccia shell interattiva
  \end{itemize}
    
  \vspace{0.25cm}

  \projectitem{UniverCity}{Completato}{Piattaforma collaborativa premiata per studenti universitari}{null}{https://github.com/davidedc97/UniverCity}
    \item Creato un sistema intelligente di condivisione appunti con algoritmi di indicizzazione avanzati, vincendo la competizione Google Workshop
    \item Presentato il progetto presso gli uffici Google di Tel Aviv, ricevendo riconoscimenti per innovazione e applicazione pratica
  \end{itemize}


\end{minipage}

\vfill
\begin{center}
  \small\textit{Per visualizzare il mio biglietto da visita digitale, apri il terminale e digita: \texttt{npx 1vcian} ✨}
\end{center}

\end{document}